\begin{table*}[t]
\centering
\small
\caption{Attack taxonomy used in the Agent-Sentinel benchmark. The evaluation spans five attack families, three tool scenarios, and five policy configurations, yielding a total of 75 scenarios.}
\label{tab:attack_taxonomy}

\begin{tabular}{llll}
\hline
\textbf{Attack Family} & \textbf{Goal} & \textbf{Representative Scenario} & \textbf{Target Surface} \\
\hline

Prompt Injection &
Manipulate agent behavior through adversarial instructions &
Injected content attempts to coerce the agent into invoking unsafe tools &
Prompt / tool boundary \\

Tool Misuse &
Abuse legitimate tools for unintended actions &
Benign tools invoked with attacker-controlled arguments &
Tool arguments \\

Capability Escalation &
Trigger operations beyond agent privileges &
Agent attempts privileged operations not allowed by policy &
Capability policy layer \\

Unauthorized API Calls &
Access disallowed services or exfiltrate data &
Outbound requests to non-allowlisted endpoints &
Network/API tools \\

Malicious System Operations &
Access or modify sensitive host resources &
Restricted file reads or shell-like operations &
Filesystem / host tools \\

\hline

\multicolumn{4}{l}{\textbf{Tool scenarios:} filesystem tools, API/network tools, system-operation workflows} \\

\multicolumn{4}{l}{\textbf{Policy configurations:} no protection, static allowlist, argument allowlist, validator-only, full Agent-Sentinel} \\

\multicolumn{4}{l}{\textbf{Total benchmark size:} $5 \times 3 \times 5 = 75$ scenarios} \\

\hline
\end{tabular}

\end{table*}
